\section{Resistor}
An ideal resistor is a two-port, passive, linear and time-invariant
device that opposes the flow of current.

\begin{itemize}
  \item \textbf{Resistance:}
    This act of opposition is called resistance and is measured in
    Ohms (\(\Omega\)). A good way to imagine resistance is to think
    about its value as a proportionality factor between voltage and
    current. In fact, This is precisely the relationship that gives
    us Ohm's law:
    \[
      V = I_ \cdot R
    \]
  \item \textbf{Conductance:}
    The inverse of resistance is called conductance (\(G\)) and is
    measured in Siemens (S). The relationship between resistance and
    conductance is:
    \[
      G = \frac{1}{R}
    \]
    We'll mostly use conductance \(G\) when we analyse resistors
    mathematically as it simplifies the analysis a lot. The rearranged
    Ohm's law in terms of conductance is:
    \[
      I = V \cdot G
    \]
  \item \textbf{Power Dissipation:}
    A resistor dissipates energy in the form of heat. The power
    dissipated by a resistor is given by:
    \[
      P = I^2 \cdot R = V^2 \cdot G
    \]
  \item \textbf{Non-Ideal Behavior:}
    In practise, resistor's are not ideal and thus most of our
    assumptions above (linear, time-invariant, dissipative) fall
    apart at very high frequencies. However for the most part,
    we can treat resistors as ideal devices.
  \item \textbf{Circuit Analysis:}
    The behaviour of resistor is not affected by the kind of Input
    Signal (AC or DC) applied to it. Therefore, The same mathematical
    model can be used to describe the behaviour of a resistor in both
    AC and DC inputs.
\end{itemize}
